\chapter{解析}
\section{背景：物理モデルを計算できる形へ移す}
解析では、実際の現象を場・ベクトル・方程式として表し、解の性質や近似の誤差を調べる。流れの保存則には発散、局所的な循環には回転、温度やポテンシャルの変化には勾配を用いる。座標変換は形の複雑な領域を計算しやすくし、行列は連立方程式や変換をまとめる。数値解析では連続な微分・積分を有限個の値から計算する。

物理専攻の読者には数学そのものより、\emph{何を入力して何を返す演算か、条件はどこで使うか}の確認が重要である。勾配はベクトル、勾配の長さと発散はスカラー、三次元の回転はベクトルである。微分前に点を代入しない。変換の向きと行・列の順序を、問題の式に合わせて読む。

\section{内積・外積と直交射影}
\label{theory:analysis:projection}
正規直交基底では$a\cdot b=\sum_i a_i b_i$であり、内積はスカラーとなる。仕事$F\cdot\Delta x$は変位方向の力成分を取り出す例である。三次元の外積$a\times b$は両者に垂直で、向きは右手系で定まり、$a\times b=-b\times a$である。平行な場合はゼロだが、それだけを根拠に交換的とはいえない。

$a$を非零の$b$に平行な$p$と垂直な$q$に分けるとき、$p=tb$、$q=a-tb$とおく。$q\cdot b=0$から
\[
 t=\frac{a\cdot b}{b\cdot b},\qquad
 p=\frac{a\cdot b}{b\cdot b}b,\qquad q=a-p
\]
となる。$b$が単位ベクトルなら分母は1である。$q\cdot b=0$だけで候補を選ぶと、$a-q$が$b$の倍数でない候補も残る。二条件を同時に満たす必要がある。

原点を通る直線への垂足の\emph{位置}は$p$であり、点$a$から垂足へ向かう\emph{変位}は$p-a$である。始点と終点を逆にすると符号が反転する。これは、点から面・線への最短距離や力の成分分解にも使える。

\section{勾配と最急上昇率}
\label{theory:analysis:gradient}
スカラー場$f$の小さな変化は$df=\nabla f\cdot dx$で表せる。単位方向$e$への方向微分は$\nabla f\cdot e$であり、コーシー・シュワルツの不等式から
\[
 \nabla f\cdot e\le |\nabla f|\,|e|=|\nabla f|
\]
となる。$\nabla f\neq0$なら$e=\nabla f/|\nabla f|$で等号が成立し、勾配方向が最急上昇、勾配のノルムがその上昇率である。勾配と方向微分の対応は大学講義の整理とも一致する\parencite{mit_chain2007}。

二変数なら$\nabla f=(f_x,f_y)$、$|\nabla f|=\sqrt{f_x^2+f_y^2}$である。求める量が「長さ」なら成分の組を答えず平方根まで取る。$\nabla f=0$の点では一次の方向微分が全方向で0であり、極値の性質には二階微分などが必要となる。

\section{発散と回転：保存と循環}
\label{theory:analysis:divcurl}
右手系の直交座標でベクトル場$V=(V_x,V_y,V_z)$に対し
\[
 \nabla\cdot V=\frac{\partial V_x}{\partial x}
 +\frac{\partial V_y}{\partial y}+\frac{\partial V_z}{\partial z},
\]
\[
 \nabla\times V=
 \left(\frac{\partial V_z}{\partial y}-\frac{\partial V_y}{\partial z},
 \frac{\partial V_x}{\partial z}-\frac{\partial V_z}{\partial x},
 \frac{\partial V_y}{\partial x}-\frac{\partial V_x}{\partial y}\right).
\]
発散は成分と同じ軸の微分を\emph{足す}。回転は別の軸の微分を\emph{引く}。これらの出力の型と座標依存性は材料工学の数学講義でも明示される\parencite{mit_vectorcalc2005}。上式は直交するデカルト座標の成分式であり、円柱・球座標へ無変更で使わない。

発散定理により、小体積から外向きに出る流束の総量を体積で割った極限が発散になる。流速を$u$とすると質量保存は
\[
 \partial_t\rho+\nabla\cdot(\rho u)=0.
\]
密度が流れに沿って不変な非圧縮性条件では$\nabla\cdot u=0$となる。「平面の全ての点で成立」という設問では、式がある曲線上でだけゼロになる候補を除く。

二次元の$V=(u,v)$を三次元へ$(u,v,0)$と埋め込むと回転の面外成分は$\partial_xv-\partial_yu$となる。過去問がこのスカラーをrotと定義した場合は、その定義に従う。三次元での回転ベクトルや発散と取り違えない。

\section{連鎖律：ヤコビ行列の向きを式から決める}
\label{theory:analysis:chain}
$x=x(r,s),y=y(r,s)$と$f=f(x,y)$に対し、全微分へ代入すると
\[
 df=f_x(x_rdr+x_sds)+f_y(y_rdr+y_sds)
\]
である。$dr,ds$の係数を比べれば
\[
 \begin{pmatrix}f_r\\f_s\end{pmatrix}
 =\begin{pmatrix}x_r&y_r\\x_s&y_s\end{pmatrix}
 \begin{pmatrix}f_x\\f_y\end{pmatrix}.
\]
一方、$f=f(x,y),g=g(x,y)$を変数とする$u(f,g)$なら
\[
 \begin{pmatrix}u_x\\u_y\end{pmatrix}
 =\begin{pmatrix}f_x&g_x\\f_y&g_y\end{pmatrix}
 \begin{pmatrix}u_f\\u_g\end{pmatrix}.
\]
ヤコビ行列という語だけから配置を決めず、各行がどの微分を返すかを確認する。連鎖律は微分記号を分数として無条件に約分する操作ではない\parencite{mit_chain2007}。

関係式$f'=1+f^2$と$f(0)$から高階導関数を求める場合も、連鎖律と積の微分を順に使う。解$f$を全区間で求めることと、一点の導関数を求めることは別の仕事である。

\section{行列式と面積要素の変数変換}
\label{theory:analysis:jacobian}
局所的に微分可能で正則な変換$(u,v)\mapsto(x,y)$では、小長方形が平行四辺形へ写り、その面積倍率は行列式の\emph{絶対値}になる。従って
\[
 dxdy=\left|\det\frac{\partial(x,y)}{\partial(u,v)}\right|dudv.
\]
向き付きの行列式と、負にならない面積の倍率を分ける。MITの変数変換講義でも面積で絶対値が必要な理由をこの幾何から説明する\parencite{mit_jacobian2007}。

例えば$x=u+v,y=uv$の微分行列は$\left(\begin{smallmatrix}1&1\\v&u\end{smallmatrix}\right)$で、行列式は$u-v$、面積倍率は$|u-v|$である。また$u,v$を交換すると同じ$x,y$になるため、積分の置換には一対一となる枝を選ぶか重複を数える処理が必要である。平成28年度の設問では符号付き行列式の候補はあるが、面積要素へ用いる符号・領域条件が不足するため疑義を保持する。

累次積分では領域も先に読む。$0\le x\le1,0\le y\le\sqrt{1-x^2}$は四分の一円である。被積分関数が$x$なら、まず$y$について積分して$\int_0^1x\sqrt{1-x^2}\,dx$とできる。被積分関数1の領域面積と、$x$の重みを持つ一次モーメントは異なる。

\section{逆行列：積で確かめ、三角構造を使う}
\label{theory:analysis:inverse}
逆行列は$AA^{-1}=A^{-1}A=I$を満たす行列である。小さい行列の随伴行列による計算と、一般の連立方程式としての意味は大学講義に整理される\parencite{mit_inverse2007}。2行2列なら
\[
 \begin{pmatrix}a&b\\c&d\end{pmatrix}^{-1}
 =\frac{1}{ad-bc}\begin{pmatrix}d&-b\\-c&a\end{pmatrix},
 \qquad ad-bc\neq0.
\]
非対角成分の符号だけでなく、分母の符号も確認する。

単位下三角行列では逆も同じ形とおき、積の非対角成分を順に0へ合わせれば
\[
 \begin{pmatrix}1&0&0\\a&1&0\\b&c&1\end{pmatrix}^{-1}
 =\begin{pmatrix}1&0&0\\-a&1&0\\ac-b&-c&1\end{pmatrix}.
\]
第3行第1列には二段の結合項$ac$が現れる。上三角の問題でも同じように、掛け算の全経路を足す。元の非対角成分を一律に符号反転するだけでは逆にならない。

実際の大規模な線形方程式$Ax=b$では、逆行列を明示的に作る代わりに分解と代入で$x$を求めることが多い。ここでの小行列の計算は写像と検算を理解するために位置付ける。

\section{数値微分：テイラー展開で次数と誤差を読む}
\label{theory:analysis:difference}
等間隔の$x_i$、幅$h>0$に対して、滑らかな関数の展開は
\[
 f_{i\pm1}=f_i\pm hf_i'+\frac{h^2}{2}f_i''
 \pm\frac{h^3}{6}f_i'''+\frac{h^4}{24}f_i^{(4)}+\cdots.
\]
差を取ると一階中心差分、和から中央を引くと二階中心差分となる。
\[
 f_i'=\frac{f_{i+1}-f_{i-1}}{2h}+O(h^2),\qquad
 f_i''=\frac{f_{i+1}-2f_i+f_{i-1}}{h^2}+O(h^2).
\]
前進・後退の一階差分は通常$O(h)$、三点の片側一階差分$(3f_i-4f_{i-1}+f_{i-2})/(2h)$は$O(h^2)$である。名称を知らなくても定数・一次・二次関数を入れれば、多くの誤った候補を除ける。二階微分の分母は長さの二乗であり、次元も検算になる。

幅を小さくすると打切り誤差は減るが、近い値の差を小さい幅で割るため丸め誤差の影響は強くなり得る。数値微分の「正しい差分式」は滑らかな関数に対する近似であり、有限の幅で任意の関数へ恒等的に正確という意味ではない。

\section{数値積分：補間とモーメント条件}
\label{theory:analysis:quadrature}
求積式を$\sum_jw_jf(x_j)$とすると、多項式$1,x,x^2,\ldots$へ正確となるための条件は
\[
 \sum_jw_jx_j^k=\int_a^b x^k\,dx
\]
である。定数関数には重み和が区間幅となる必要がある。対称区間では奇関数部分は相殺されるため、偶数次数の条件を確認しやすい。これは式の係数を暗記するだけでなく、似た誤答式を検証する方法となる。

$a,b,c$が等間隔$h$の三点なら、二次補間多項式を積分するシンプソン法は
\[
 \int_a^cf(x)\,dx\approx\frac h3[f(a)+4f(b)+f(c)].
\]
全幅を$H=c-a=2h$と書けば係数は$H/6$である。記号の定義を取り違えない。滑らかな関数の一パネル誤差は$-h^5f^{(4)}(\xi)/90$であり、三次以下の多項式には正確となる\parencite{nist_quadrature}。

等間隔に縛られず、区間$[-1,1]$の二点を$\pm t$、等しい重みを$w$とすると、定数と二次のモーメント条件から$2w=2$、$2wt^2=2/3$となる。従って$w=1,t=1/\sqrt3$である。奇数次数は対称性で相殺されるので三次以下に正確な二点ガウス求積
\[
 \int_{-1}^1f(x)\,dx=f(-1/\sqrt3)+f(1/\sqrt3)
 \quad (\deg f\le3)
\]
が得られる。評価点が正しくても重みが違えば別の式である。

近似手法が指定された問題では、その手法で計算した値を答える。例えば$f(x)=1/(x+3)$を$[-1,1]$で一パネルのシンプソン法にかけると$25/36\approx0.69444$であり、厳密値$\ln2\approx0.693147$と違う。両者に最も近い候補が同じでも、導出を置き換えない。

\section{固有振動：復元力のモデル化}
\subsection{背景}
支持条件や接続を変えると、同じ材質・同じ部品でも振動特性は変わる。まず変位に対する復元力を求め、慣性との組合せを考える。
\subsection{理論と重要式}
一自由度の線形系では$m\ddot{x}+k_{\rm eff}x=0$から$\omega_0^2=k_{\rm eff}/m$。
質量が同じなら、有効剛性だけで固有振動数の大小を比較できる。
直列では同じ力を受けるばねの変形が加算される。並列では同じ変位を受けるばねの力が加算される。
\subsection{技術的意味}
左右の固定壁につながる質点を右へ変位させると、左のばねの伸びと右のばねの圧縮はともに左向きの復元力を生む。
図の左右という見かけから力を相殺してしまわない。
\subsection{一般の多自由度系への接続}
$M\ddot{q}+Kq=0$から$(K-\omega^2M)\phi=0$が得られる。
ただし過去問の主題がばね・はりの振動比較なら主分類は振動・力学とし、一般化固有値問題は横断参照とする。


\section{数値解析：誤差を発生箇所で分ける}
\label{theory:numerics:errors}
物理モデルを計算する際には、材料や境界条件を単純化するモデル誤差、連続な場を格子や多項式で表す離散化誤差、方程式を途中で解くのを止める反復誤差、有限の桁で演算する丸め誤差がある。実験と計算の差をすべてメッシュの粗さへ帰してはいけない。同じ評価量を固定し、どの操作で結果が変わるかを切り分ける。

例えば形状を変えず格子幅を半分にして応力が変わるなら離散化の影響が疑われる。格子を固定して停止閾値を厳しくすると値が変わるなら反復が不足している。倍精度へ変える操作は仮数の精度を増すが、モデルの誤りや切り捨てた高次項を復元しない。反復回数が多いから閾値を緩めるという操作は、速度と引換えに許容誤差を増す操作である。

\subsection{有限精度：桁落ちと情報落ち}
\label{theory:numerics:floating}
浮動小数点数は指数と有限桁の仮数で値を表す。大きな数へ小さい数を足すと、指数を揃えた小さい側の桁が保持範囲外へ落ちることがある。近い二数の差では共通の大きい部分が消え、各入力に既に含まれた誤差が差の値に対して大きくなる。これは異なる現象である。正確に表された二数の減算が常に新しい大誤差を作るわけではない\parencite{goldberg1991}。

実数の結合則は浮動小数点演算の同じ結果を保証しない。通常の倍精度で$10^{16}+1$の1が失われる場合、$(10^{16}+1)-10^{16}$と$(10^{16}-10^{16})+1$は異なる。桁落ちを避ける式変形の例は
\[
 \sqrt{x+1}-\sqrt{x}=\frac{1}{\sqrt{x+1}+\sqrt{x}}\quad(x\ge0)
\]
である。大きい$x$で左辺を近い数の差として直接計算するより、右辺は差による増幅を避けられる。ただしオーバーフローや入力精度など別の限界は残る。

\subsection{差分式：整合性と精度を別に調べる}
\label{theory:numerics:diff}
前節\ref{theory:analysis:difference}のTaylor展開から、同じ一階導関数を近似しても前進・後退差分の誤差は$O(h)$、中心差分は$O(h^2)$となる。二次精度の後退差分も
\[
 \frac{3f_i-4f_{i-1}+f_{i-2}}{2h}=f'_i+O(h^2)
\]
と導ける。定数項の係数和は$3-4+1=0$、一次項は$2hf'_i$、二次項は消える。補間多項式の微分から差分と余項を得る整理もある\parencite{nist_differentiation}。

式が目的の導関数に近づくことが整合性であり、誤差が減る次数が精度である。$(f_{i+1}-2f_i+f_{i-1})/h$は$hf_i''+O(h^3)$であって、一階微分に整合しない。分母を$h^2$に直しても二階微分である。検算には$f=1,x,x^2$を代入するとよい。格子を過度に細かくすると近い値の差の丸め誤差が増し、打切り誤差の低下だけから最終精度を断定できない。

\subsection{補間：与えられた点と元の関数を区別する}
\label{theory:numerics:interpolation}
相異なる$n+1$点を通る次数$n$以下の多項式は一意である。二つあれば差の多項式は$n+1$個の根を持つが次数は$n$以下であるため、差は零となる。Lagrange形式では
\[
 p(x)=\sum_{j=0}^{n}f(x_j)\prod_{k\ne j}\frac{x-x_k}{x_j-x_k}.
\]
各基底は自分の節点で1、他の節点で0となる\parencite{nist_interpolation}。三点なら$p=ax^2+bx+c$へ代入して解く方法も使える。補間は与えられた点を正確に通すが、元の関数全体を正確に復元するとは限らない。区間外の外挿では、補間値の一致だけで信頼性を判断しない。

\subsection{積分領域：三角形の境界を先に書く}
\label{theory:numerics:integration}
頂点$(0,0),(a,0),(0,a)$の三角形は$0\le x\le a,\ 0\le y\le a-x$である。長方形$[0,a]^2$へ置き換えると領域を増やす。一次関数の例では
\[
 \int_0^a\int_0^{a-x}(x+y)\,dy\,dx
 =\int_0^a\left[x(a-x)+\frac{(a-x)^2}{2}\right]dx
 =\frac{a^3}{3}.
\]
独立の幾何検算として、一次関数の三角形内の平均値は重心での値$2a/3$に等しく、面積$a^2/2$との積も$a^3/3$となる。単位も長さの三乗で確認できる。

\subsection{求積：重み和と奇偶の検算}
\label{theory:numerics:quadrature}
前節\ref{theory:analysis:quadrature}のモーメント条件を使うと、$[-1,1]$の一次関数$ax+b$に正確な公式は重み和2かつ$\sum_jw_jx_j=0$を満たす。端点の台形則、中央点則、三点のSimpson則はいずれも満たす。台形則の一般関数に対する余項には二階導関数が現れる\parencite{nist_quadrature}。

例えば端点重みを$1/4$ずつ、中央重みを1とすると重み和は$3/2$であり、定数$b$を$3b/2$へ誤る。$b=0$なら偶然ゼロへ一致するが、全ての一次関数に正確な一般式ではない。特定の値で一致することと、ある関数族に正確な公式であることを分ける。

\subsection{Newton法：更新の符号と終了条件}
\label{theory:numerics:newton}
現在値$x_n$で接線を引き、その零点を次の値とする。$f(x_n)+f'(x_n)(x_{n+1}-x_n)=0$から
\[
 x_{n+1}=x_n-\frac{f(x_n)}{f'(x_n)}.
\]
分母は導関数であり積ではない。単根近傍の条件が満たされると二次収束するが、初期値を問わず収束する保証はなく、$f'(x_n)=0$ではこの更新を使えない\parencite{nist_newton}。

フローチャートが$x\leftarrow x-\Delta x$なら$\Delta x=f/f'$、$x\leftarrow x+\Delta x$なら$\Delta x=-f/f'$である。小さいステップで終了する分岐には$|\Delta x|<\varepsilon$を合わせる。実際の計算ではステップが小さくても残差が十分小さいとは限らないため、$|f(x)|$、相対ステップ、最大反復数、導関数の小ささを別に確認する。試験の簡略図に合わせることと、実務の停止条件が十分であることを分ける。

\section{有限要素法：場を節点値で近似する}
\label{theory:fem:background}
連続体の変位を有限個の節点値で表し、材料則とつり合いを離散方程式へ移す。変位を$u_h=\sum_iN_i u_i$、微分演算を$L$とすると、ひずみは$\varepsilon=Lu_h=Bq$、線形弾性の応力は$\sigma=D\varepsilon$である。形状関数の微分を含む$B$を用い、要素剛性は$K_e=\int_{V_e}B^TDB\,dV$となる。この補間から剛性行列への関係はMIT講義でも整理される\parencite{mit_roy_fea1999}。以下の具体的な形状関数と幾何計算は独自に導出する。

\subsection{一次三角形：一定ひずみの理由}
\label{theory:fem:triangle}
非退化の三頂点を持つ一次三角形では、各変位成分を$u=a+bx+cy$と表す。微分は$u_x=b,u_y=c$であるため、微小ひずみは要素内で一定となる。均質な線形弾性なら応力も一定である。6節点の二次三角形には$x^2,xy,y^2$の項が加わり、微分は一般に一次関数となる。

これは任意の節点変位に対する要素の性質である。二次三角形でも二次項がゼロとなる特別な場ならひずみは一定になり得る。「一定ひずみ要素」を選ぶ設問で、その特殊例だけを理由に高次要素まで選ばない。退化三角形や不均質な材料では、一定ひずみから一定応力への移行にも条件が要る。

\subsection{双一次四角形：交差項が残る}
\label{theory:fem:q4}
自然座標の四隅を$(-1,-1),(1,-1),(1,1),(-1,1)$とすると
\[
 (N_1,N_2,N_3,N_4)=\frac14\bigl((1-\xi)(1-\eta),(1+\xi)(1-\eta),
 (1+\xi)(1+\eta),(1-\xi)(1+\eta)\bigr).
\]
各節点で対応する関数だけ1となり、和は1である。展開した$\xi,\eta,\xi\eta$の係数ベクトルは順に
\[
 \tfrac14(-1,1,1,-1),\quad\tfrac14(-1,-1,1,1),\quad
 \tfrac14(1,-1,1,-1).
\]
非定数の各係数の和が0という検算もできる。

長方形の物理座標なら$u=a+bx+cy+dxy$となり、$u_x=b+dy,u_y=c+dx$には位置依存が残る。一般四角形では座標写像のJacobianを介して微分し、ひずみはさらに形状に依存する。四隅4節点というだけで要素内一定応力と判断しない。

\subsection{変位の連続と応力の跳び}
\label{theory:fem:continuity}
適合する変位法では隣接要素が同じ辺の節点変位を共有する。一次の辺補間が同じなので辺上の変位も一致する。しかし法線方向の微分までは一般に一致せず、微分から作るひずみや応力に跳びが生じ得る。一様ひずみの場では隣接応力が一致する例もあるため、「常に不連続」という全称は強すぎる。

表示用に節点応力を平均すると滑らかな色分布になることがあるが、平均処理が近似精度の保証を与えるわけではない。変位、積分点応力、平均応力のどれを比較したかを明示する。

\subsection{メッシュ：局所応力に上下保証はない}
\label{theory:fem:mesh}
細分化するh改良は空間解像度を増し、高次化するp改良は近似関数の表現力を増す。滑らかな解、健全な形状、適切な積分や境界条件の下で有効となる。高勾配域へ要素を配置する考え方はMITの例とも対応する\parencite{mit_roy_fea1999}。

線形弾性の適合変位法における全体のエネルギーの傾向から、任意の点の応力が上界であるとはいえない。粗いメッシュでは応力集中を過小評価することもある。「変形が小さいので応力は常に高く安全側」という結論は成立しない。異なる荷重では重要な勾配の位置も変わり、一つの荷重で妥当だった分割を無検証で使い回さない。

複数のメッシュで同じ変位、反力、位置の応力を比較する。尖った角や点荷重に理想化した特異点ではピーク応力が細分ごとに増す場合があり、有限値への収束を当然視しない。要素の歪みや退化、丸め誤差、材料モデルの誤差は別に確認する。

\subsection{面積座標：重心・内心・外心を分ける}
\label{theory:fem:area}
三角形ABC内の点Pに対し、$S_A=S_{PBC},S_B=S_{PCA},S_C=S_{PAB}$とおき、$\lambda_i=S_i/S$と定める。非負で和は1となる。頂点Aでは$(1,0,0)$、重心では$(1/3,1/3,1/3)$である。面積座標は一次三角形の形状関数にもなる。

内心は三辺からの距離$r$が等しく、$S_A=r|BC|/2$などとなる。従って座標は向かい辺の長さの比である。$AB:BC:CA=3:4:5$なら$(\lambda_A,\lambda_B,\lambda_C)=(4,5,3)/12$となる。Bが直角で、外心は斜辺ACの中点なので$(1/2,0,1/2)$である。重心を内心や外心と混同せず、図の見た目より辺長比を優先する。三角形外の点に拡張する場合には符号付き面積が必要となる。

\subsection{二層の熱伝導：並列と直列の加算量}
\label{theory:fem:thermal}
定常伝導の平板では$Q=kA\Delta T/d$であり、熱流は温度勾配に比例する\parencite{openstax_thermal_conduction}。接触熱抵抗がなく、厚さ$l_1,l_2$の層に平行な流れでは温度勾配が共通であり、各断面の熱流を足す。層を横切る流れでは熱流が共通であり、各層の温度差を足す。$L=l_1+l_2$として
\[
 k_\parallel=\frac{k_1l_1+k_2l_2}{L},\qquad
 k_\perp=\frac{L}{l_1/k_1+l_2/k_2}.
\]
前者は厚さで重み付けした算術平均、後者は調和平均である。正の熱伝導率では両者とも各材料値の間に入り、$k_\parallel\ge k_\perp$となる。界面熱抵抗や過渡現象があれば、この簡単な二式のモデルをそのまま適用しない。

厚さを材料に逆に対応させた誤答式は等厚の特殊値で正しい式と一致する。試験では一般式を導く目的と読みつつ、その一致と時間条件の省略を解説の疑義へ残す。

\section{軸変形は釣合と適合を同時に満たす}
\label{theory:analysis:fixedbar}
\label{theory:analysis:compositebar}
\label{theory:analysis:parallelbar}
\label{theory:analysis:extension}
一様な棒の軸剛性は \(EA/l\)、伸びは \(Nl/(EA)\) である。力の釣合だけで決まらない反力には、固定端間の総変位がゼロという適合条件を加える。途中の荷重点を動かすと、左の伸びと右の縮みは同じ変位となり、両区間は荷重点に対して復元剛性を並列に提供する。右向きの外力に対する支持反力を右向き正で表せば、いずれも負になる。

二本の等長棒を回転しない剛体端板で結ぶ場合は伸び、従ってひずみが共通である。応力比は弾性係数比、軸力比は \(EA\) の比となる。端板が回転する条件や長さが違う場合はこの結論をそのまま使わない。計算は N と mm なら \(E\) を N/mm\(^2\) にそろえる。\cite{mit_material_elastic2000}

\section{横収縮と非圧縮極限}
\label{theory:analysis:hooke}
\label{theory:analysis:incompressible}
\label{theory:analysis:thermal}
等方線形弾性体の構成則は
\[
 \varepsilon_x=\frac{\sigma_x-\nu(\sigma_y+\sigma_z)}{E},\qquad
 \varepsilon_v=\frac{1-2\nu}{E}(\sigma_x+\sigma_y+\sigma_z).
\]
引張の横収縮は負符号で入る。任意の応力状態でも体積を変えない極限は \(\nu\to1/2\) である。このとき \(K=E/[3(1-2\nu)]\) は発散し、通常の有限圧縮性の逆構成則へそのまま代入する操作とは区別する。熱を受ける完全固定棒では全ひずみ \(\sigma/E+\alpha\Delta T=0\) から温度上昇に対する圧縮応力を得る。部分拘束では支持の変形も加えて解く。\cite{mit_material_elastic2000}

\section{弾性エネルギーの基準と安定性}
\label{theory:analysis:axialenergy}
\label{theory:analysis:potential}
\label{theory:analysis:verticalenergy}
線形ばねの蓄積仕事は力変位線図の三角形面積である。
\[
 U=\frac12 P\delta=\frac{P^2l}{2EA},\qquad
 U_s=\frac12ku^2.
\]
ばね自身の \(u\) は自然長から測る。鉛直ばねで平衡からの変位を \(x\) とすると \(u=mg/k+x\) であり、ばねエネルギーと重力エネルギーを合計して平衡値を引いた振動エネルギーは \(kx^2/2\) となる。ばねだけの全量とは違う。全ポテンシャル \(\Pi=ku^2/2-mgu\) は平衡で停留し、二階微分が正なので安定な最小となる。\cite{analysis57_shm}

\section{曲げ方向の高さと片持ち撓み}
\label{theory:analysis:section}
\label{theory:analysis:cantilever}
矩形断面の曲げ軸に関する断面二次モーメントは \(I=bh^3/12\)。断面積が等しくても高さ \(h\) が違えば曲げ剛性は違う。同じ \(E\) と荷重に対し小さい \(I\) が大きい撓みを与える。先端荷重を受ける細長い片持ち梁では
\[
 \delta=\frac{PL^3}{3EI}.
\]
他の支持や分布荷重なら係数も式も変わる。寸法比を求めるときは各梁の長さと断面を別々に代入し、無次元比で比較する。\cite{struct_mit_beam,analysis57_beam}

\section{孔縁の自由表面にも接線応力が残る}
\label{theory:analysis:inglis}
自由表面は応力テンソル全体がゼロという条件ではなく、法線 \(n\) に対する表面力 \(\sigma n=0\) である。楕円孔の水平端では法線が水平なので法線垂直応力とせん断応力が消え、接線垂直応力は残る。十分広い等方線形弾性板で遠方荷重が楕円長軸に直交する場合、孔端の応力集中係数は \(1+2a/b\)。円孔の3へ戻ることを検算に使える。有限幅、塑性や鋭い亀裂先端には別のモデルが必要となる。\cite{analysis57_inglis}

\section{小振動の剛性を配置から組み立てる}
\label{theory:analysis:spring}
\label{theory:analysis:springgravity}
\label{theory:analysis:stringmass}
運動を平衡からの変位で表すと \(m\ddot x+k_{\rm eff}x=0\) なので \(\omega=\sqrt{k_{\rm eff}/m}\)、\(f=\omega/(2\pi)\)。ばね並列は剛性の和、直列はコンプライアンスの和になる。質点を左右の壁へつなぐ二ばねは、一方が伸び他方が縮んでも復元力が同じ方向へ働く。重力の一定成分は線形ばねの平衡位置をずらすが剛性を変えない。

張った弦の中央に質点を置く場合は横変位に対する張力成分を小角近似で加える。各半長を \(L\) とすると \(k_{\rm eff}=2T/L\)。全長と半長を取り違えず、質量の変更も二乗振動数の比へ含める。\cite{analysis57_shm}

\section{梁振動と柱座屈の境界条件}
\label{theory:analysis:beamfreq}
\label{theory:analysis:buckling}
一様Euler--Bernoulli梁の曲げ固有角振動数は \(\omega_1=\beta_1^2\sqrt{EI/(\rho A)}/L^2\)。片持ちの特性式は \(\cosh\beta\cos\beta=-1\)、両端固定は \(=1\) で、それぞれ非零の最小根は約1.875と4.730である。単純支持では \(\beta_1=\pi\)。剛体モードのゼロ根を曲げ基本振動と混同しない。\cite{analysis57_beam}

Euler座屈は \(P_{\rm cr}=\pi^2EI/(KL)^2\) に支持を有効長係数 \(K\) として含める。ピン・ピンは1、固定・自由は2、固定・ピンは約0.699。後者の特性根 \(z\) は \(\tan z=z\)、\(K=\pi/z\) である。これは真直な細長柱の弾性理想モデルで、初期不整や短柱圧壊の設計耐力を保証しない。

\section{節点の力から伸びへ進む}
\label{theory:analysis:truss}
\label{theory:analysis:trussforce}
ピン節点の二力部材は軸方向だけに力を伝える。節点の水平・鉛直力を釣り合わせて部材力 \(N_i\) を求め、その後に \(\delta_i=N_il_i/(EA_i)\) を使う。軸力比と伸び比は長さ・断面の因子が違う。角度が水平基準か鉛直基準かを先に定め、正負によって引張・圧縮を分ける。\cite{mit_material_elastic2000}

\section{回転慣性と動力の二つの関係}
\label{theory:analysis:inertia}
\label{theory:analysis:torque}
\label{theory:analysis:geartorque}
\label{theory:analysis:pendulum}
回転運動は \(\tau=I\alpha\)、回転動力は \(P=\tau\omega\)。角速度を時間で変える問題と減速比の問題では、同じ式だけでは解けない。入力速度が出力の \(n\) 倍、効率1なら入力トルクは出力の \(1/n\)。換算慣性 \(I/n^2\) を使うときの角加速度も入力側へ換算する。

一様薄板の面外重心軸には垂直軸定理 \(I_z=I_x+I_y\) が使える。面外に厚さがある物体には同じ式を無条件に使わない。物理振子の微小角運動は \(I\ddot\theta+Mgd\theta=0\)、従って周期は \(2\pi\sqrt{I/(Mgd)}\)。支点から重心までの距離 \(d\) と物体の慣性モーメントを別々に求める。

\section{終端・定常の意味を保存則で読む}
\label{theory:analysis:terminal}
\label{theory:analysis:bernoulli}
\label{theory:analysis:motional}
線形抵抗落下は \(m\dot v=mg-kv\)。終端では \(\dot v=0\) で、速度依存が二乗抵抗とは違う。水平な定常・非圧縮・非粘性流れには \(A_av_a=A_bv_b\) と \(p+\rho v^2/2=\text{一定}\) を併用する。断面積比が速度二乗へ入ること、粘性損失や高さの差を原典条件から外せるかを確認する。

運動起電力の棒では \(\mathcal E=BLv\)、\(I=\mathcal E/R\)、制動力 \(IBL\)。定常で駆動力と制動が釣り合う。抵抗が大きいほど電流が小さくなるため制動が弱まるという極限も検算に使える。

\section{抵抗網は線の形でなく節点の接続で解く}
\label{theory:analysis:network}
\label{theory:analysis:tetra}
同じ二節点を結ぶ枝は並列、途中で分岐のない経路が直列である。交差線の接続点を確認し、対称性で同電位と証明できる節点をまとめる。任意網には内部節点で
\[
 \sum_j\frac{V_i-V_j}{R_{ij}}=0
\]
を解き、端子電圧を1とすれば \(R_{\rm eq}=1/I_{\rm terminal}\)。一辺の中点を端子にする場合はその辺を半抵抗二本に分ける。四面体の頂点・対辺中点・対辺の両中点は異なる端子条件であり、同じ正四面体だから同抵抗とはならない。\cite{analysis57_resistors}

\section{拡散の減衰と幾何の重心}
\label{theory:analysis:heat}
\label{theory:analysis:centroid}
正の拡散係数の通常の熱伝導は、固定温度又は断熱の安定境界下で空間固有値 \(\lambda\ge0\) に対し \(e^{-\kappa\lambda t}\) の応答を持つ。時間一階微分という形式だけでは非振動性を証明できず、連立一階方程式が回転運動を表す例もある。微分次数、演算子、境界条件を一組で扱う。

面積重心は \(\bar x=A^{-1}\int x\,dA\)、\(\bar y=A^{-1}\int y\,dA\)。一様面、線、体積の重心では積分の測度が違う。対称性で二座標が等しい場合も、値は面積一次モーメントから求める。四分円面と円弧を同じ式で扱わない。

\subsection{平成23年度 問3-1}
\label{q:h23-reg-basic-3-1}

\noindent\textbf{原典：} sources/H23/basic.pdf、PDF物理8頁。

\textbf{公式正答：} 2（single）。

\par\textit{原典本文の校正と独自解説は未作成。分類のみの登録項目。}

\subsection{平成23年度 問3-2}
\label{q:h23-reg-basic-3-2}

\noindent\textbf{原典：} sources/H23/basic.pdf、PDF物理9頁。

\textbf{公式正答：} 1（single）。

\par\textit{原典本文の校正と独自解説は未作成。分類のみの登録項目。}

\subsection{平成23年度 問3-3}
\label{q:h23-reg-basic-3-3}

\noindent\textbf{原典：} sources/H23/basic.pdf、PDF物理10頁。

\textbf{公式正答：} 5（single）。

\par\textit{原典本文の校正と独自解説は未作成。分類のみの登録項目。}

\subsection{平成23年度 問3-4}
\label{q:h23-reg-basic-3-4}

\noindent\textbf{原典：} sources/H23/basic.pdf、PDF物理11頁。

\textbf{公式正答：} 3（single）。

\par\textit{原典本文の校正と独自解説は未作成。分類のみの登録項目。}

\subsection{平成23年度 問3-5}
\label{q:h23-reg-basic-3-5}

\noindent\textbf{原典：} sources/H23/basic.pdf、PDF物理11頁。

\textbf{公式正答：} 1（single）。

\par\textit{原典本文の校正と独自解説は未作成。分類のみの登録項目。}

\subsection{平成24年度 問3-1}
\label{q:h24-reg-basic-3-1}

\noindent\textbf{原典：} sources/H24/basic.pdf、PDF物理8頁。

\textbf{公式正答：} 1（single）。

\par\textit{原典本文の校正と独自解説は未作成。分類のみの登録項目。}

\subsection{平成24年度 問3-2}
\label{q:h24-reg-basic-3-2}

\noindent\textbf{原典：} sources/H24/basic.pdf、PDF物理8頁。

\textbf{公式正答：} 2（single）。

\paragraph{問題の要旨} ばね接続から有効剛性を求め、3系の固有振動数を比較する。

\paragraph{必要知識} ばねの直列・並列合成、微小変位の復元力、固有角振動数

\paragraph{考え方（査読待ち初稿）} aでは復元力が2kxなので有効剛性2k。bでは変形量が加算され1/keff=1/k+1/k、よってkeff=k/2。cも左右の復元力が同方向に合計され2kx。質量が同じため固有角振動数はsqrt(keff/m)で比較でき、bのみ最小。正答②。

\paragraph{他の選択肢（査読待ち初稿）} ①aと③cはいずれも有効剛性2k。④は最小でないaを、⑤は最小でないcを含む。cの左右の力は打ち消し合わず、変位に対して同じ復元方向に働く。

\paragraph{関連問題} \hyperref[q:h26-reg-basic-3-5]{H26-REG-BASIC-3-5}

\par 難易度 2（校準初稿）、優先度未評価。

\subsection{平成24年度 問3-3}
\label{q:h24-reg-basic-3-3}

\noindent\textbf{原典：} sources/H24/basic.pdf、PDF物理9頁。

\textbf{公式正答：} 1（single）。

\par\textit{原典本文の校正と独自解説は未作成。分類のみの登録項目。}

\subsection{平成24年度 問3-4}
\label{q:h24-reg-basic-3-4}

\noindent\textbf{原典：} sources/H24/basic.pdf、PDF物理10頁。

\textbf{公式正答：} 4（single）。

\par\textit{原典本文の校正と独自解説は未作成。分類のみの登録項目。}

\subsection{平成24年度 問3-5}
\label{q:h24-reg-basic-3-5}

\noindent\textbf{原典：} sources/H24/basic.pdf、PDF物理11頁。

\textbf{公式正答：} 4（single）。

\par\textit{原典本文の校正と独自解説は未作成。分類のみの登録項目。}

\subsection{平成25年度 問3-1}
\label{q:h25-reg-basic-3-1}

\noindent\textbf{原典：} sources/H25/basic.pdf、PDF物理10頁。

\textbf{公式正答：} 5（single）。

\par\textit{原典本文の校正と独自解説は未作成。分類のみの登録項目。}

\subsection{平成25年度 問3-2}
\label{q:h25-reg-basic-3-2}

\noindent\textbf{原典：} sources/H25/basic.pdf、PDF物理11頁。

\textbf{公式正答：} 3（single）。

\par\textit{原典本文の校正と独自解説は未作成。分類のみの登録項目。}

\subsection{平成25年度 問3-3}
\label{q:h25-reg-basic-3-3}

\noindent\textbf{原典：} sources/H25/basic.pdf、PDF物理12頁。

\textbf{公式正答：} 1（single）。

\par\textit{原典本文の校正と独自解説は未作成。分類のみの登録項目。}

\subsection{平成25年度 問3-4}
\label{q:h25-reg-basic-3-4}

\noindent\textbf{原典：} sources/H25/basic.pdf、PDF物理13頁。

\textbf{公式正答：} 5（single）。

\par\textit{原典本文の校正と独自解説は未作成。分類のみの登録項目。}

\subsection{平成25年度 問3-5}
\label{q:h25-reg-basic-3-5}

\noindent\textbf{原典：} sources/H25/basic.pdf、PDF物理13頁。

\textbf{公式正答：} 2（single）。

\par\textit{原典本文の校正と独自解説は未作成。分類のみの登録項目。}

\subsection{平成25年度 問3-6}
\label{q:h25-reg-basic-3-6}

\noindent\textbf{原典：} sources/H25/basic.pdf、PDF物理14頁。

\textbf{公式正答：} 4（single）。

\par\textit{原典本文の校正と独自解説は未作成。分類のみの登録項目。}

\subsection{平成26年度 問3-1}
\label{q:h26-reg-basic-3-1}

\noindent\textbf{原典：} sources/H26/basic.pdf、PDF物理9頁。

\textbf{公式正答：} 2（single）。

\par\textit{原典本文の校正と独自解説は未作成。分類のみの登録項目。}

\subsection{平成26年度 問3-2}
\label{q:h26-reg-basic-3-2}

\noindent\textbf{原典：} sources/H26/basic.pdf、PDF物理9頁。

\textbf{公式正答：} 3（single）。

\par\textit{原典本文の校正と独自解説は未作成。分類のみの登録項目。}

\subsection{平成26年度 問3-3}
\label{q:h26-reg-basic-3-3}

\noindent\textbf{原典：} sources/H26/basic.pdf、PDF物理10頁。

\textbf{公式正答：} 5（single）。

\par\textit{原典本文の校正と独自解説は未作成。分類のみの登録項目。}

\subsection{平成26年度 問3-4}
\label{q:h26-reg-basic-3-4}

\noindent\textbf{原典：} sources/H26/basic.pdf、PDF物理10頁。

\textbf{公式正答：} 1（single）。

\par\textit{原典本文の校正と独自解説は未作成。分類のみの登録項目。}

\subsection{平成26年度 問3-5}
\label{q:h26-reg-basic-3-5}

\noindent\textbf{原典：} sources/H26/basic.pdf、PDF物理11頁。

\textbf{公式正答：} 2（single）。

\par\textit{原典本文の校正と独自解説は未作成。分類のみの登録項目。}

\subsection{平成26年度 問3-6}
\label{q:h26-reg-basic-3-6}

\noindent\textbf{原典：} sources/H26/basic.pdf、PDF物理12頁。

\textbf{公式正答：} 2（single）。

\par\textit{原典本文の校正と独自解説は未作成。分類のみの登録項目。}

\subsection{平成27年度 問3-1}
\label{q:h27-reg-basic-3-1}

\noindent\textbf{原典：} sources/H27/basic.pdf、PDF物理8頁。

\textbf{公式正答：} 2（single）。

\par\textit{原典本文の校正と独自解説は未作成。分類のみの登録項目。}

\subsection{平成27年度 問3-2}
\label{q:h27-reg-basic-3-2}

\noindent\textbf{原典：} sources/H27/basic.pdf、PDF物理8頁。

\textbf{公式正答：} 2（single）。

\par\textit{原典本文の校正と独自解説は未作成。分類のみの登録項目。}

\subsection{平成27年度 問3-3}
\label{q:h27-reg-basic-3-3}

\noindent\textbf{原典：} sources/H27/basic.pdf、PDF物理9頁。

\textbf{公式正答：} 1（single）。

\par\textit{原典本文の校正と独自解説は未作成。分類のみの登録項目。}

\subsection{平成27年度 問3-4}
\label{q:h27-reg-basic-3-4}

\noindent\textbf{原典：} sources/H27/basic.pdf、PDF物理9頁。

\textbf{公式正答：} 3（single）。

\par\textit{原典本文の校正と独自解説は未作成。分類のみの登録項目。}

\subsection{平成27年度 問3-5}
\label{q:h27-reg-basic-3-5}

\noindent\textbf{原典：} sources/H27/basic.pdf、PDF物理10頁。

\textbf{公式正答：} 3（single）。

\par\textit{原典本文の校正と独自解説は未作成。分類のみの登録項目。}

\subsection{平成27年度 問3-6}
\label{q:h27-reg-basic-3-6}

\noindent\textbf{原典：} sources/H27/basic.pdf、PDF物理11頁。

\textbf{公式正答：} 4（single）。

\par\textit{原典本文の校正と独自解説は未作成。分類のみの登録項目。}

\subsection{平成28年度 問3-1}
\label{q:h28-reg-basic-3-1}

\noindent\textbf{原典：} sources/H28/basic.pdf、PDF物理12頁。

\textbf{公式正答：} 3（single）。

\par\textit{原典本文の校正と独自解説は未作成。分類のみの登録項目。}

\subsection{平成28年度 問3-2}
\label{q:h28-reg-basic-3-2}

\noindent\textbf{原典：} sources/H28/basic.pdf、PDF物理12頁。

\textbf{公式正答：} 3（single）。

\par\textit{原典本文の校正と独自解説は未作成。分類のみの登録項目。}

\subsection{平成28年度 問3-3}
\label{q:h28-reg-basic-3-3}

\noindent\textbf{原典：} sources/H28/basic.pdf、PDF物理13頁。

\textbf{公式正答：} 2（single）。

\par\textit{原典本文の校正と独自解説は未作成。分類のみの登録項目。}

\subsection{平成28年度 問3-4}
\label{q:h28-reg-basic-3-4}

\noindent\textbf{原典：} sources/H28/basic.pdf、PDF物理14頁。

\textbf{公式正答：} 5（single）。

\par\textit{原典本文の校正と独自解説は未作成。分類のみの登録項目。}

\subsection{平成28年度 問3-5}
\label{q:h28-reg-basic-3-5}

\noindent\textbf{原典：} sources/H28/basic.pdf、PDF物理15頁。

\textbf{公式正答：} 1（single）。

\par\textit{原典本文の校正と独自解説は未作成。分類のみの登録項目。}

\subsection{平成28年度 問3-6}
\label{q:h28-reg-basic-3-6}

\noindent\textbf{原典：} sources/H28/basic.pdf、PDF物理16頁。

\textbf{公式正答：} 4（single）。

\par\textit{原典本文の校正と独自解説は未作成。分類のみの登録項目。}

\subsection{平成29年度 問3-1}
\label{q:h29-reg-basic-3-1}

\noindent\textbf{原典：} sources/H29/basic.pdf、PDF物理11頁。

\textbf{公式正答：} 5（single）。

\par\textit{原典本文の校正と独自解説は未作成。分類のみの登録項目。}

\subsection{平成29年度 問3-2}
\label{q:h29-reg-basic-3-2}

\noindent\textbf{原典：} sources/H29/basic.pdf、PDF物理11頁。

\textbf{公式正答：} 4（single）。

\par\textit{原典本文の校正と独自解説は未作成。分類のみの登録項目。}

\subsection{平成29年度 問3-3}
\label{q:h29-reg-basic-3-3}

\noindent\textbf{原典：} sources/H29/basic.pdf、PDF物理11頁。

\textbf{公式正答：} 4（single）。

\par\textit{原典本文の校正と独自解説は未作成。分類のみの登録項目。}

\subsection{平成29年度 問3-4}
\label{q:h29-reg-basic-3-4}

\noindent\textbf{原典：} sources/H29/basic.pdf、PDF物理12頁。

\textbf{公式正答：} 2（single）。

\par\textit{原典本文の校正と独自解説は未作成。分類のみの登録項目。}

\subsection{平成29年度 問3-5}
\label{q:h29-reg-basic-3-5}

\noindent\textbf{原典：} sources/H29/basic.pdf、PDF物理13頁。

\textbf{公式正答：} 3（single）。

\par\textit{原典本文の校正と独自解説は未作成。分類のみの登録項目。}

\subsection{平成29年度 問3-6}
\label{q:h29-reg-basic-3-6}

\noindent\textbf{原典：} sources/H29/basic.pdf、PDF物理14頁。

\textbf{公式正答：} 5（single）。

\par\textit{原典本文の校正と独自解説は未作成。分類のみの登録項目。}

\subsection{平成30年度 問3-1}
\label{q:h30-reg-basic-3-1}

\noindent\textbf{原典：} sources/H30/basic.pdf、PDF物理11頁。

\textbf{公式正答：} 1（single）。

\par\textit{原典本文の校正と独自解説は未作成。分類のみの登録項目。}

\subsection{平成30年度 問3-2}
\label{q:h30-reg-basic-3-2}

\noindent\textbf{原典：} sources/H30/basic.pdf、PDF物理11頁。

\textbf{公式正答：} 4（single）。

\par\textit{原典本文の校正と独自解説は未作成。分類のみの登録項目。}

\subsection{平成30年度 問3-3}
\label{q:h30-reg-basic-3-3}

\noindent\textbf{原典：} sources/H30/basic.pdf、PDF物理12頁。

\textbf{公式正答：} 2（single）。

\par\textit{原典本文の校正と独自解説は未作成。分類のみの登録項目。}

\subsection{平成30年度 問3-4}
\label{q:h30-reg-basic-3-4}

\noindent\textbf{原典：} sources/H30/basic.pdf、PDF物理13頁。

\textbf{公式正答：} 2（single）。

\par\textit{原典本文の校正と独自解説は未作成。分類のみの登録項目。}

\subsection{平成30年度 問3-5}
\label{q:h30-reg-basic-3-5}

\noindent\textbf{原典：} sources/H30/basic.pdf、PDF物理14頁。

\textbf{公式正答：} 3（single）。

\par\textit{原典本文の校正と独自解説は未作成。分類のみの登録項目。}

\subsection{平成30年度 問3-6}
\label{q:h30-reg-basic-3-6}

\noindent\textbf{原典：} sources/H30/basic.pdf、PDF物理14頁。

\textbf{公式正答：} 3（single）。

\par\textit{原典本文の校正と独自解説は未作成。分類のみの登録項目。}

\subsection{令和元年度 問3-1}
\label{q:r01-reg-basic-3-1}

\noindent\textbf{原典：} sources/R01/basic.pdf、PDF物理12頁。

\textbf{公式正答：} 5（single）。

\par\textit{原典本文の校正と独自解説は未作成。分類のみの登録項目。}

\subsection{令和元年度 問3-2}
\label{q:r01-reg-basic-3-2}

\noindent\textbf{原典：} sources/R01/basic.pdf、PDF物理12頁。

\textbf{公式正答：} 5（single）。

\par\textit{原典本文の校正と独自解説は未作成。分類のみの登録項目。}

\subsection{令和元年度 問3-3}
\label{q:r01-reg-basic-3-3}

\noindent\textbf{原典：} sources/R01/basic.pdf、PDF物理13頁。

\textbf{公式正答：} 1（single）。

\par\textit{原典本文の校正と独自解説は未作成。分類のみの登録項目。}

\subsection{令和元年度 問3-4}
\label{q:r01-reg-basic-3-4}

\noindent\textbf{原典：} sources/R01/basic.pdf、PDF物理13頁。

\textbf{公式正答：} 4（single）。

\par\textit{原典本文の校正と独自解説は未作成。分類のみの登録項目。}

\subsection{令和元年度 問3-5}
\label{q:r01-reg-basic-3-5}

\noindent\textbf{原典：} sources/R01/basic.pdf、PDF物理14頁。

\textbf{公式正答：} 4（single）。

\par\textit{原典本文の校正と独自解説は未作成。分類のみの登録項目。}

\subsection{令和元年度 問3-6}
\label{q:r01-reg-basic-3-6}

\noindent\textbf{原典：} sources/R01/basic.pdf、PDF物理15頁。

\textbf{公式正答：} 3（single）。

\par\textit{原典本文の校正と独自解説は未作成。分類のみの登録項目。}

\subsection{令和元年度再試験 問3-1}
\label{q:r01-ret-basic-3-1}

\noindent\textbf{原典：} sources/R01\_retest/basic.pdf、PDF物理11頁。

\textbf{公式正答：} 5（single）。

\par\textit{原典本文の校正と独自解説は未作成。分類のみの登録項目。}

\subsection{令和元年度再試験 問3-2}
\label{q:r01-ret-basic-3-2}

\noindent\textbf{原典：} sources/R01\_retest/basic.pdf、PDF物理11頁。

\textbf{公式正答：} 2（single）。

\par\textit{原典本文の校正と独自解説は未作成。分類のみの登録項目。}

\subsection{令和元年度再試験 問3-3}
\label{q:r01-ret-basic-3-3}

\noindent\textbf{原典：} sources/R01\_retest/basic.pdf、PDF物理11頁。

\textbf{公式正答：} 5（single）。

\par\textit{原典本文の校正と独自解説は未作成。分類のみの登録項目。}

\subsection{令和元年度再試験 問3-4}
\label{q:r01-ret-basic-3-4}

\noindent\textbf{原典：} sources/R01\_retest/basic.pdf、PDF物理12頁。

\textbf{公式正答：} 5（single）。

\par\textit{原典本文の校正と独自解説は未作成。分類のみの登録項目。}

\subsection{令和元年度再試験 問3-5}
\label{q:r01-ret-basic-3-5}

\noindent\textbf{原典：} sources/R01\_retest/basic.pdf、PDF物理12頁。

\textbf{公式正答：} 4（single）。

\par\textit{原典本文の校正と独自解説は未作成。分類のみの登録項目。}

\subsection{令和元年度再試験 問3-6}
\label{q:r01-ret-basic-3-6}

\noindent\textbf{原典：} sources/R01\_retest/basic.pdf、PDF物理13頁。

\textbf{公式正答：} 2（single）。

\par\textit{原典本文の校正と独自解説は未作成。分類のみの登録項目。}

\subsection{令和2年度 問3-1}
\label{q:r02-reg-basic-3-1}

\noindent\textbf{原典：} sources/R02/basic.pdf、PDF物理11頁。

\textbf{公式正答：} 2（single）。

\par\textit{原典本文の校正と独自解説は未作成。分類のみの登録項目。}

\subsection{令和2年度 問3-2}
\label{q:r02-reg-basic-3-2}

\noindent\textbf{原典：} sources/R02/basic.pdf、PDF物理11頁。

\textbf{公式正答：} 4（single）。

\par\textit{原典本文の校正と独自解説は未作成。分類のみの登録項目。}

\subsection{令和2年度 問3-3}
\label{q:r02-reg-basic-3-3}

\noindent\textbf{原典：} sources/R02/basic.pdf、PDF物理11頁。

\textbf{公式正答：} 3（single）。

\par\textit{原典本文の校正と独自解説は未作成。分類のみの登録項目。}

\subsection{令和2年度 問3-4}
\label{q:r02-reg-basic-3-4}

\noindent\textbf{原典：} sources/R02/basic.pdf、PDF物理12頁。

\textbf{公式正答：} 4（single）。

\par\textit{原典本文の校正と独自解説は未作成。分類のみの登録項目。}

\subsection{令和2年度 問3-5}
\label{q:r02-reg-basic-3-5}

\noindent\textbf{原典：} sources/R02/basic.pdf、PDF物理13頁。

\textbf{公式正答：} 1（single）。

\par\textit{原典本文の校正と独自解説は未作成。分類のみの登録項目。}

\subsection{令和2年度 問3-6}
\label{q:r02-reg-basic-3-6}

\noindent\textbf{原典：} sources/R02/basic.pdf、PDF物理14頁。

\textbf{公式正答：} 5（single）。

\par\textit{原典本文の校正と独自解説は未作成。分類のみの登録項目。}

\subsection{令和3年度 問3-1}
\label{q:r03-reg-basic-3-1}

\noindent\textbf{原典：} sources/R03/basic.pdf、PDF物理11頁。

\textbf{公式正答：} 4（single）。

\par\textit{原典本文の校正と独自解説は未作成。分類のみの登録項目。}

\subsection{令和3年度 問3-2}
\label{q:r03-reg-basic-3-2}

\noindent\textbf{原典：} sources/R03/basic.pdf、PDF物理11頁。

\textbf{公式正答：} 2（single）。

\par\textit{原典本文の校正と独自解説は未作成。分類のみの登録項目。}

\subsection{令和3年度 問3-3}
\label{q:r03-reg-basic-3-3}

\noindent\textbf{原典：} sources/R03/basic.pdf、PDF物理12頁。

\textbf{公式正答：} 1（single）。

\par\textit{原典本文の校正と独自解説は未作成。分類のみの登録項目。}

\subsection{令和3年度 問3-4}
\label{q:r03-reg-basic-3-4}

\noindent\textbf{原典：} sources/R03/basic.pdf、PDF物理13頁。

\textbf{公式正答：} 3（single）。

\par\textit{原典本文の校正と独自解説は未作成。分類のみの登録項目。}

\subsection{令和3年度 問3-5}
\label{q:r03-reg-basic-3-5}

\noindent\textbf{原典：} sources/R03/basic.pdf、PDF物理14頁。

\textbf{公式正答：} 4（single）。

\par\textit{原典本文の校正と独自解説は未作成。分類のみの登録項目。}

\subsection{令和3年度 問3-6}
\label{q:r03-reg-basic-3-6}

\noindent\textbf{原典：} sources/R03/basic.pdf、PDF物理15頁。

\textbf{公式正答：} 5（single）。

\par\textit{原典本文の校正と独自解説は未作成。分類のみの登録項目。}

\subsection{令和4年度 問3-1}
\label{q:r04-reg-basic-3-1}

\noindent\textbf{原典：} sources/R04/basic.pdf、PDF物理10頁。

\textbf{公式正答：} 4（single）。

\par\textit{原典本文の校正と独自解説は未作成。分類のみの登録項目。}

\subsection{令和4年度 問3-2}
\label{q:r04-reg-basic-3-2}

\noindent\textbf{原典：} sources/R04/basic.pdf、PDF物理10頁。

\textbf{公式正答：} 2（single）。

\par\textit{原典本文の校正と独自解説は未作成。分類のみの登録項目。}

\subsection{令和4年度 問3-3}
\label{q:r04-reg-basic-3-3}

\noindent\textbf{原典：} sources/R04/basic.pdf、PDF物理11頁。

\textbf{公式正答：} 4（single）。

\par\textit{原典本文の校正と独自解説は未作成。分類のみの登録項目。}

\subsection{令和4年度 問3-4}
\label{q:r04-reg-basic-3-4}

\noindent\textbf{原典：} sources/R04/basic.pdf、PDF物理11頁。

\textbf{公式正答：} 2（single）。

\par\textit{原典本文の校正と独自解説は未作成。分類のみの登録項目。}

\subsection{令和4年度 問3-5}
\label{q:r04-reg-basic-3-5}

\noindent\textbf{原典：} sources/R04/basic.pdf、PDF物理12頁。

\textbf{公式正答：} 3（single）。

\par\textit{原典本文の校正と独自解説は未作成。分類のみの登録項目。}

\subsection{令和4年度 問3-6}
\label{q:r04-reg-basic-3-6}

\noindent\textbf{原典：} sources/R04/basic.pdf、PDF物理13頁。

\textbf{公式正答：} 5（single）。

\par\textit{原典本文の校正と独自解説は未作成。分類のみの登録項目。}

\subsection{令和5年度 問3-1}
\label{q:r05-reg-basic-3-1}

\noindent\textbf{原典：} sources/R05/basic.pdf、PDF物理9頁。

\textbf{公式正答：} 4（single）。

\par\textit{原典本文の校正と独自解説は未作成。分類のみの登録項目。}

\subsection{令和5年度 問3-2}
\label{q:r05-reg-basic-3-2}

\noindent\textbf{原典：} sources/R05/basic.pdf、PDF物理9頁。

\textbf{公式正答：} 2（single）。

\par\textit{原典本文の校正と独自解説は未作成。分類のみの登録項目。}

\subsection{令和5年度 問3-3}
\label{q:r05-reg-basic-3-3}

\noindent\textbf{原典：} sources/R05/basic.pdf、PDF物理10頁。

\textbf{公式正答：} 1（single）。

\par\textit{原典本文の校正と独自解説は未作成。分類のみの登録項目。}

\subsection{令和5年度 問3-4}
\label{q:r05-reg-basic-3-4}

\noindent\textbf{原典：} sources/R05/basic.pdf、PDF物理10頁。

\textbf{公式正答：} 5（single）。

\par\textit{原典本文の校正と独自解説は未作成。分類のみの登録項目。}

\subsection{令和5年度 問3-5}
\label{q:r05-reg-basic-3-5}

\noindent\textbf{原典：} sources/R05/basic.pdf、PDF物理11頁。

\textbf{公式正答：} 2（single）。

\par\textit{原典本文の校正と独自解説は未作成。分類のみの登録項目。}

\subsection{令和5年度 問3-6}
\label{q:r05-reg-basic-3-6}

\noindent\textbf{原典：} sources/R05/basic.pdf、PDF物理12頁。

\textbf{公式正答：} 3（single）。

\par\textit{原典本文の校正と独自解説は未作成。分類のみの登録項目。}

\subsection{令和6年度 問3-1}
\label{q:r06-reg-basic-3-1}

\noindent\textbf{原典：} sources/R06/basic.pdf、PDF物理10頁。

\textbf{公式正答：} 4（single）。

\par\textit{原典本文の校正と独自解説は未作成。分類のみの登録項目。}

\subsection{令和6年度 問3-2}
\label{q:r06-reg-basic-3-2}

\noindent\textbf{原典：} sources/R06/basic.pdf、PDF物理10頁。

\textbf{公式正答：} 2（single）。

\par\textit{原典本文の校正と独自解説は未作成。分類のみの登録項目。}

\subsection{令和6年度 問3-3}
\label{q:r06-reg-basic-3-3}

\noindent\textbf{原典：} sources/R06/basic.pdf、PDF物理11頁。

\textbf{公式正答：} 5（single）。

\par\textit{原典本文の校正と独自解説は未作成。分類のみの登録項目。}

\subsection{令和6年度 問3-4}
\label{q:r06-reg-basic-3-4}

\noindent\textbf{原典：} sources/R06/basic.pdf、PDF物理12頁。

\textbf{公式正答：} 3（single）。

\par\textit{原典本文の校正と独自解説は未作成。分類のみの登録項目。}

\subsection{令和6年度 問3-5}
\label{q:r06-reg-basic-3-5}

\noindent\textbf{原典：} sources/R06/basic.pdf、PDF物理13頁。

\textbf{公式正答：} 3（single）。

\par\textit{原典本文の校正と独自解説は未作成。分類のみの登録項目。}

\subsection{令和6年度 問3-6}
\label{q:r06-reg-basic-3-6}

\noindent\textbf{原典：} sources/R06/basic.pdf、PDF物理14頁。

\textbf{公式正答：} 3（single）。

\par\textit{原典本文の校正と独自解説は未作成。分類のみの登録項目。}

\subsection{令和7年度 問3-1}
\label{q:r07-reg-basic-3-1}

\noindent\textbf{原典：} sources/R07/basic.pdf、PDF物理8頁。

\textbf{公式正答：} 2（single）。

\par\textit{原典本文の校正と独自解説は未作成。分類のみの登録項目。}

\subsection{令和7年度 問3-2}
\label{q:r07-reg-basic-3-2}

\noindent\textbf{原典：} sources/R07/basic.pdf、PDF物理8頁。

\textbf{公式正答：} 1（single）。

\par\textit{原典本文の校正と独自解説は未作成。分類のみの登録項目。}

\subsection{令和7年度 問3-3}
\label{q:r07-reg-basic-3-3}

\noindent\textbf{原典：} sources/R07/basic.pdf、PDF物理9頁。

\textbf{公式正答：} 5（single）。

\par\textit{原典本文の校正と独自解説は未作成。分類のみの登録項目。}

\subsection{令和7年度 問3-4}
\label{q:r07-reg-basic-3-4}

\noindent\textbf{原典：} sources/R07/basic.pdf、PDF物理10頁。

\textbf{公式正答：} 2（single）。

\par\textit{原典本文の校正と独自解説は未作成。分類のみの登録項目。}

\subsection{令和7年度 問3-5}
\label{q:r07-reg-basic-3-5}

\noindent\textbf{原典：} sources/R07/basic.pdf、PDF物理11頁。

\textbf{公式正答：} 5（single）。

\par\textit{原典本文の校正と独自解説は未作成。分類のみの登録項目。}

\subsection{令和7年度 問3-6}
\label{q:r07-reg-basic-3-6}

\noindent\textbf{原典：} sources/R07/basic.pdf、PDF物理12頁。

\textbf{公式正答：} 1（single）。

\par\textit{原典本文の校正と独自解説は未作成。分類のみの登録項目。}

